\section{Численный алгоритм}
\label{sec:algorithm}

\subsection{Прямая задача}
\label{sec:forward-method}

Прямой расчёт начинается с определения стационарной температуры для начальной конфигурации, после чего последовательно рассматриваются все эпохи. На каждом временном шаге решается задача теплопереноса, а по полученному изменению температуры вычисляются приращения аррениусовских интегралов и отражательная способность витринита. Переход между эпохами сопровождается переносом полей на следующую сетку, как показано в алгоритме~\ref{alg:forward}.

\bigskip
\noindent\begin{minipage}{\linewidth}
\hrule
\medskip
\refstepcounter{reportalgorithm}\label{alg:forward}
\textbf{Алгоритм~\thereportalgorithm. Решение прямой задачи}

\smallskip
\textbf{Вход.} Конфигурации, свойства среды, граничные данные, кинетические параметры модели EASY\%Ro и параметры численного метода.

\textbf{Результат.} Поля температуры и отражательной способности витринита.

\begin{enumerate}
\setlength{\itemsep}{0pt}
\setlength{\parsep}{0pt}
\item Построить сетки эпох и задать движение ячеек по заданной геометрии осадочного бассейна.
\item Найти начальную температуру как решение стационарной задачи методом Пикара.
\item Обнулить аррениусовские интегралы и задать моменты отложения материала.
\item \textbf{Для каждой эпохи выполнить}
\item \hspace*{1em}\textbf{Пока не достигнут конец эпохи, повторять}
\item \hspace*{2em}Выбрать шаг времени, не выходящий за конец эпохи.
\item \hspace*{2em}Обновить геометрию, пористость и граничные данные.
\item \hspace*{2em}Взять температуру прошлого шага начальным приближением.
\item \hspace*{2em}\textbf{Пока не достигнута сходимость метода Пикара, повторять}
\item \hspace*{3em}Обновить свойства среды и собрать линейную систему.
\item \hspace*{3em}Найти поправку методом GMRES с предобуславливателем ILU(2).
\item \hspace*{3em}Обновить температуру и вычислить нелинейную невязку.
\item \hspace*{2em}\textbf{Конец цикла Пикара}
\item \hspace*{2em}Обновить аррениусовские интегралы и отражательную способность витринита.
\item \hspace*{2em}Перейти к новому моменту времени.
\item \hspace*{1em}\textbf{Конец цикла по времени}
\item \hspace*{1em}При наличии следующей эпохи перенести температуру и аррениусовские интегралы на её сетку.
\item \textbf{Конец цикла по эпохам}
\item Вернуть рассчитанные поля.
\end{enumerate}
\smallskip
\hrule
\end{minipage}

\subsubsection{Расчётная сетка}
\label{sec:mesh}

Расчётная область разбивается на вертикальные колонки прямоугольных ячеек. Внутри эпохи горизонтальные координаты колонок и число ячеек остаются постоянными, а вертикальные координаты границ ячеек изменяются в соответствии с движением материала. Горизонтальный и вертикальный шаги $\Delta x,\Delta z$ задаются независимо. Это позволяет выбирать более мелкое разбиение по глубине, где свойства среды меняются при переходе между слоями, и сохранять более крупный шаг вдоль вытянутого разреза.

В пределах каждой колонки кровля и подошва геологического слоя совпадают с горизонтальными границами расчётных ячеек. При этом уровни границ в соседних колонках могут различаться, поэтому вертикальная сторона одной ячейки может соприкасаться с несколькими ячейками другой колонки. Каждый участок такого соприкосновения называется вертикальной подгранью и рассматривается отдельно при аппроксимации теплового потока. Неизвестная температура $T_P$ относится к центру ячейки $P$. Устройство сетки и взаимное смещение соседних колонок показаны на рис.~\ref{fig:mesh}.

\reportfigure{02_mesh.png} {Расчётная сетка из вертикальных колонок для Модели~1 в современном состоянии. (а) сетка всего разреза с отмеченными участками увеличения; (б) участок вблизи разлома; (в) участок при $600\leq x\leq900$~м и $2400\leq z\leq3200$~м, на котором видны несовпадающие уровни ячеек соседних колонок.} {fig:mesh}

\subsubsection{Пространственный оператор и граничные потоки}

Для пространственной аппроксимации уравнения теплопереноса применяется метод конечных объёмов~\cite{samarskii2003}. Уравнение~\eqref{eq:heat} интегрируется по контрольному объёму, совпадающему с движущейся ячейкой $P(t)$, площадь которой обозначена через $|P|$. Длина её грани обозначена через $|\sigma|$, а внешняя единичная нормаль к грани~--- через $\boldsymbol n_P$. Поскольку ячейка перемещается вместе с материалом, сумма временного и конвективного слагаемых аппроксимируется полной производной температуры ячейки по времени. Объёмная теплоёмкость $C_P=(\rho c)_P$ и объёмная плотность радиогенного тепловыделения $A_P$ принимаются равными значениям в центре ячейки. Для интегралов временного слагаемого и источника используются значения в центре ячейки, а теорема о дивергенции преобразует пространственный оператор в сумму потоков через грани

\begin{gather}
 \int_{P(t)} \rho c\left(\frac{\partial T}{\partial t}
              + \boldsymbol v \cdot \nabla T\right)\,d\Omega
      \approx C_P|P|\frac{dT_P}{dt},\qquad
 \int_{P(t)} A\,d\Omega \approx A_P|P|,\\
 \int_{P(t)} \nabla \cdot (\lambda\nabla T)\,d\Omega
       = -\sum_{\sigma\subset\partial P} F_{P\sigma},\qquad
 F_{P\sigma} = -\int_\sigma \lambda\nabla T \cdot \boldsymbol n_P\,d\sigma.
\end{gather}

Поток через внутреннюю грань входит в уравнения двух соседних ячеек с противоположными знаками, $F_{N\sigma}=-F_{P\sigma}$. Благодаря этому при суммировании уравнений обмен теплом между ячейками взаимно сокращается. На горизонтальных гранях центры соседних ячеек лежат на одной нормали, что позволяет применять двухточечную аппроксимацию. Расстояния от центров до грани обозначаются через $d_{P\sigma}$ и $d_{N\sigma}$. Сопротивления участков между центрами и гранью складываются последовательно, поэтому интегральный поток имеет вид

\begin{equation}
 F_{P\sigma} = \frac{|\sigma|(T_P - T_N)}
                    {d_{P\sigma}/\lambda_P + d_{N\sigma}/\lambda_N}.
\end{equation}

На вертикальной подграни направление между центрами соседних ячеек в общем случае не совпадает с нормалью, поэтому разность двух температур в центрах ячеек содержит вклад изменения температуры как по горизонтали, так и по вертикали и не определяет нормальную производную. Для учёта этого смещения применяется многоточечная аппроксимация потока (MPFA). Помимо центральных температур в ней используются значения на соседних горизонтальных границах ячеек. Соотношение направлений между центрами и нормалей к граням показано на рис.~\ref{fig:flux}.

С каждой стороны подграни строится одностороннее приближение плотности потока, направленного из соответствующей ячейки, и полученные значения обозначаются через $f_{P\sigma}$ и $f_{N\sigma}$. Обозначим $g_P=\lambda_P/d_{P\sigma}$, $g_N=\lambda_N/d_{N\sigma}$, где расстояния до вертикальной подграни измерены по горизонтали. Со стороны $P$ вводится вспомогательное значение
\begin{gather}
 \omega = \frac{g_N}{g_P + g_N},\qquad
 T_Q = (1 - \omega)T_P + \omega T_N,\qquad
 z_Q = z_P + \omega(z_N - z_P),\\
 f_{P\sigma} = \frac{\lambda_P}{d_{P\sigma}}
 \left[(T_P - T_Q) - \frac{z_Q - z_P}{z_V - z_P}(T_P - T_V)\right].
 \label{eq:one-sided-flux}
\end{gather}
Точка $Q$ расположена на вертикали подграни, а точка $V$ выбирается на верхней или нижней горизонтальной границе ячейки $P$ с противоположной стороны по глубине относительно $Q$. Значение $T_V$ определяется по тепловым сопротивлениям соседних ячеек либо по граничному условию. Второе слагаемое в~\eqref{eq:one-sided-flux} учитывает вертикальный сдвиг центров. Со стороны $N$ построение выполняется аналогично. В линейном варианте односторонние потоки согласуются с весами, зависящими от нормальных проводимостей
\begin{equation}
 F_{P\sigma} = |\sigma|\left(
 \frac{g_N}{g_P + g_N}f_{P\sigma}
 - \frac{g_P}{g_P + g_N}f_{N\sigma}\right).
 \label{eq:mpfa-combination}
\end{equation}
Для нелинейной многоточечной аппроксимации потока (NMPFA) веса зависят от значений односторонних потоков. В обоих случаях при фиксированных коэффициентах поток имеет вид

\begin{equation}
 F_{P\sigma} = \sum_{N \in \mathcal N_\sigma} a_{\sigma N}T_N + b_\sigma.
\end{equation}

Множество $\mathcal N_\sigma$ включает ячейки, температуры которых участвуют в аппроксимации данного потока. Коэффициенты $a_{\sigma N}$ зависят от геометрии и теплопроводности, а в NMPFA также от текущего температурного поля. Слагаемое $b_\sigma$ учитывает заданные граничные температуры и потоки. Принципы построения многоточечных аппроксимаций рассмотрены в~\cite[гл.~2]{inmost2020}. В основных расчётах используется линейный вариант MPFA, а нелинейный вариант NMPFA применяется для сопоставления. При этом линейность MPFA означает линейную зависимость потока от температур только при фиксированной теплопроводности. Полная задача остаётся нелинейной вследствие температурной зависимости свойств среды.

\reportfigure{03_flux.png} {Расположение ячеек и направлений, используемых при аппроксимации теплового потока. Из центра ячейки $P$ проведены направления к центрам соседних ячеек $N_0,N_1,N_2$. Красным выделены горизонтальная грань и две вертикальные подграни, чёрными стрелками показаны их нормали. Цвет ячеек обозначает разные теплопроводности $\lambda_1$ и $\lambda_2$.} {fig:flux}

Граничные потоки входят в интегральное уравнение теплопереноса для ячеек, прилегающих к границе области. На горизонтальном участке поверхности используется заданная температура $T_{\mathrm{surf}}$, а на нижней границе поток $q_{\mathrm{basal}}$, заданный на единицу горизонтальной проекции, умножается на ширину ячейки $\Delta x_P$. С учётом теплоизоляции боковых границ получаются соотношения

\begin{equation}
 F_{P\sigma}^{\mathrm{surf}} = \frac{\lambda_P|\sigma|}{d_{P\sigma}}
                    (T_P - T_{\mathrm{surf}}),\qquad
 F_{P\sigma}^{\mathrm{basal}} = -q_{\mathrm{basal}}\Delta x_P,\qquad
 F_{P\sigma}^{\mathrm{side}} = 0.
\end{equation}

Ступенчатое представление геометрии образует вертикальные участки поверхности, на которых заданная температура учитывается в односторонней многоточечной аппроксимации. Вертикальные участки нижней границы осадочного чехла имеют нулевую горизонтальную проекцию, поэтому приток тепла из фундамента определяется горизонтальными гранями. Суммирование потоков и учёт радиогенного тепловыделения дают систему пространственно дискретизованных уравнений

\begin{equation}
 C_P|P|\frac{dT_P}{dt} + \mathcal L_P(T) = 0,\qquad
 \mathcal L_P(T) = \sum_{\sigma \subset \partial P} F_{P\sigma}(T) - A_P|P|.
\end{equation}

\subsubsection{Интегрирование по времени}

Система интегрируется неявным методом Эйлера с шагом $\Delta t=t^{n+1}-t^n$, при этом геометрия, граничные данные и теплофизические коэффициенты определяются в момент $t^{n+1}$. Поскольку $T_P^n$ и $T_P^{n+1}$ относятся к одной движущейся ячейке, отношение разности температур к шагу времени аппроксимирует материальную производную, и уравнение временного шага принимает вид

\begin{equation}
 C_P^{n+1}|P^{n+1}|\frac{T_P^{n+1} - T_P^n}{\Delta t}
       + \mathcal L_P(T^{n+1}) = 0.
 \label{eq:euler}
\end{equation}

Неявный метод Эйлера имеет первый порядок аппроксимации по времени. Если очередной шаг выходит за конец эпохи, его длина сокращается до момента $t_{k+1}$. Таким образом, переход между эпохами и перенос полей на следующую сетку выполняются на границах временных шагов. Начальная стационарная задача решается с тем же пространственным оператором и начальными граничными данными при отсутствии временного слагаемого, то есть $\mathcal L_P(T_0)=0$.

\subsubsection{Решение нелинейной системы}

Из-за температурной зависимости теплопроводности $\lambda(T)$ и объёмной теплоёмкости $\rho c(T)$ система~\eqref{eq:euler} нелинейна и решается методом Пикара. На итерации $m$ свойства вычисляются по текущему температурному полю и фиксируются при определении поправки, в NMPFA аналогично фиксируются веса односторонних потоков. Через $\boldsymbol R$ обозначен вектор невязок уравнений временного шага, через $M(T^{(m)})$~--- матрица с фиксированными коэффициентами, а поправка и следующее приближение определяются соотношениями

\begin{equation}
 M(T^{(m)})\,\delta T = -\boldsymbol R(T^{(m)}),\qquad
 T^{(m+1)} = T^{(m)} + \delta T.
\end{equation}

Метод Пикара позволяет формировать $M$ из коэффициентов дискретного уравнения без вычисления их производных по температуре, необходимых в методе Ньютона. После нахождения поправки свойства и нелинейная невязка обновляются, а итерации завершаются по абсолютному допуску её евклидовой нормы либо по относительному допуску к начальной норме. Начальным приближением служит температурное поле предыдущего временного шага.

\subsubsection{Решение линейной системы}

Линейные системы на итерациях Пикара решаются обобщённым методом минимальных невязок GMRES~\cite{saad1986}, применимым к несимметричным матрицам, возникающим при выбранной аппроксимации потоков. Используется перезапуск после 100~итераций и предобуславливание неполным LU-разложением с уровнем заполнения 2, обозначаемым ILU(2), а сходимость определяется по норме предобусловленной невязки. Выбор сочетания GMRES(100)/ILU(2) основан на сопоставлении времени решения и температурных полей с LU-разложением в библиотеке UMFPACK, приведённом в разделе~\ref{sec:implementation}.

\subsubsection{Вычисление отражательной способности витринита}

После завершения нелинейных итераций вычисляются приращения аррениусовских интегралов, определяющих отражательную способность витринита. В каждой материальной точке температура внутри шага $[t^n,t^{n+1}]$ считается линейной по времени. Интегрирование начинается с момента $t_* = \max(t^n,t_{\mathrm{dep}})$, чтобы в температурную историю не входило время до отложения материала. Температура $T_a=T(t_*)$ определяется линейной интерполяцией во времени, а $T_b=T(t^{n+1})$ соответствует концу шага. Если $t_*\geq t^{n+1}$, приращение интеграла принимается равным нулю.

При обозначении $\beta_i=E_i/R_g$ первообразная $\Phi_i$ функции $\exp(-\beta_i/T)$ выражается через интегральную экспоненту $\operatorname{Ei}$
\begin{equation}
 G(\chi) = \frac{e^{-\chi}}{\chi} + \operatorname{Ei}(-\chi),\qquad
 \Phi_i(T) = \beta_i G(\beta_i/T).
 \label{eq:ro-primitive}
\end{equation}
При $T_a\ne T_b$ линейная зависимость температуры от времени позволяет заменить переменную интегрирования и выразить приращение через значения первообразной~\eqref{eq:ro-primitive} на концах интервала
\begin{equation}
 \Delta I_i = \frac{A_i(t^{n+1} - t_*)}{T_b - T_a}
             \bigl(\Phi_i(T_b) - \Phi_i(T_a)\bigr).
 \label{eq:ro-increment}
\end{equation}
При постоянной температуре скорость реакции также постоянна, поэтому её интеграл равен произведению скорости на длительность интервала
\begin{equation}
 \Delta I_i = A_i(t^{n+1} - t_*)e^{-\beta_i/T_a},\qquad T_a=T_b.
 \label{eq:ro-increment-constant}
\end{equation}

Приращение по формуле~\eqref{eq:ro-increment} или~\eqref{eq:ro-increment-constant} добавляется к накопленному интегралу $I_i$, после чего вычисляется отражательная способность $R_o$. Моменты отложения определяются соотношением~\eqref{eq:deposition}, а расположение точек расчёта интегралов и определение температуры в них описаны в разделе~\ref{sec:implementation}.

\subsection{Обратная задача}
\label{sec:inverse-method}

Для каждого пробного набора параметров решается прямая задача по алгоритму~\ref{alg:forward}, включая пересчёт начальной стационарной температуры, зависящей от базального потока. По современным значениям температуры и отражательной способности в точках наблюдений вычисляется вектор нормированных невязок, используемый для выбора следующих параметров в алгоритме~\ref{alg:inverse}.

\bigskip
\noindent\begin{minipage}{\linewidth}
\hrule
\medskip
\refstepcounter{reportalgorithm}\label{alg:inverse}
\textbf{Алгоритм~\thereportalgorithm. Решение обратной задачи}

\smallskip
\textbf{Вход.} Данные прямой задачи, скважинные наблюдения с оценками погрешностей, начальное приближение потока и параметры оптимизации.

\textbf{Результат.} Найденные значения параметров базального теплового потока.

\begin{enumerate}
\setlength{\itemsep}{0pt}
\setlength{\parsep}{0pt}
\item Подготовить геометрию, сетки эпох и точки наблюдений.
\item Задать допустимое начальное приближение параметров и ограничения.
\item \textbf{Пока не выполнен критерий остановки, повторять}
\item \hspace*{1em}Выбрать очередной набор параметров методом POUNDERS.
\item \hspace*{1em}Определить по параметрам историю базального теплового потока.
\item \hspace*{1em}Решить прямую задачу по алгоритму~\ref{alg:forward}.
\item \hspace*{1em}Вычислить значения полей в точках наблюдений.
\item \hspace*{1em}Вычислить нормированные невязки и целевую функцию.
\item \hspace*{1em}Сохранить лучшее найденное приближение параметров.
\item \hspace*{1em}Учесть невязки при построении локальных приближений POUNDERS.
\item \hspace*{1em}Проверить критерии остановки.
\item \textbf{Конец цикла подбора параметров}
\item Вернуть найденные параметры базального теплового потока.
\end{enumerate}
\smallskip
\hrule
\end{minipage}

\subsubsection{Параметризация потока}

Искомый поток $q(t)$ задаётся значениями $q_k=q(t_k)$ на границах эпох и линейно интерполируется между ними. В Модели~3 используются 17~узлов и столько же искомых параметров. Ограничения на диапазон и монотонность потока обеспечиваются переходом к безразмерным параметрам $\eta_k\in[0,1]$, из которых узловые значения вычисляются последовательно

\begin{equation}
 q_0 = q_{\min} + (q_{\max} - q_{\min})\eta_0,\qquad
 q_k = q_{\min} + (q_{k-1} - q_{\min})\eta_k,\quad k \geq 1.
\end{equation}

Здесь $q_{\min}$ и $q_{\max}$~--- нижняя и верхняя границы допустимого потока. При любых $\eta_k\in[0,1]$ выполняются условия $q_{\min}\leq q_0\leq q_{\max}$ и $q_{\min}\leq q_k\leq q_{k-1}$ при $k\geq1$, поэтому поток не возрастает с течением модельного времени. Значение $\eta_k=1$ при $k\geq1$ соответствует постоянному потоку в пределах эпохи.

\subsubsection{Минимизация целевой функции}

Для совместного учёта наблюдений невязки температуры и отражательной способности нормируются на оценки погрешностей $\sigma_j^T$ и $\sigma_j^R$, равные 2\% от модулей измеренных значений. Температурная погрешность определяется в градусах Цельсия, индекс $j$ нумерует наблюдения соответствующей величины во всех скважинах, а нормированные невязки задаются соотношениями
\begin{gather}
 \sigma_j^T = 0{,}02\left|T_{C,j}^{\mathrm{obs}}\right|,\qquad
 \sigma_j^R = 0{,}02\left|R_{o,j}^{\mathrm{obs}}\right|,\\
 r_j^T = \frac{T_{C,j}^{\mathrm{calc}} - T_{C,j}^{\mathrm{obs}}}{\sigma_j^T},\qquad
 r_j^{R_o} = \frac{R_{o,j}^{\mathrm{calc}} - R_{o,j}^{\mathrm{obs}}}{\sigma_j^R}.
 \label{eq:normalized-residuals}
\end{gather}
Такая нормировка соответствует весам $w_j^T = (\sigma_j^T)^{-2}$ и $w_j^R = (\sigma_j^R)^{-2}$ в функционале~\eqref{eq:objective}. Минимизируется сумма квадратов нормированных невязок
\begin{equation}
 J = J_T + J_{R_o} = \sum_j (r_j^T)^2 + \sum_j (r_j^{R_o})^2.
\end{equation}
Слагаемые $J_T$ и $J_{R_o}$ характеризуют отклонения расчётных температур и отражательной способности от измерений с учётом принятых погрешностей. Регуляризующее слагаемое не вводится.

Для минимизации используется POUNDERS~\cite{wild2017,petscManual}, предназначенный для нелинейных задач наименьших квадратов без вычисления производных по параметрам. Метод строит локальные приближения отдельных компонент вектора невязок и по ним выбирает следующий набор параметров, для которого выполняется новый прямой расчёт.

При $p$~параметрах начальное построение локальных приближений требует $p+1$~прямых расчётов для исходного и возмущённых наборов, то есть при $p=17$ до первого обновления выполняются 18~расчётов. Поэтому при оценке вычислительных затрат учитываются как итерации метода оптимизации, так и число решений прямой задачи.

\subsubsection{Условия остановки}

Условие $|r_j|\leq1$ для невязок~\eqref{eq:normalized-residuals} означает, что отклонение расчётного значения от измеренного не превышает принятой погрешности. Подбор завершается при выполнении этого условия для всех наблюдений, достижении внутреннего критерия POUNDERS либо заданного ограничения числа итераций или прямых расчётов.

В результате прямой алгоритм последовательно определяет температуру и отражательную способность на заданной геометрической истории, а обратный использует этот расчёт для подбора потока по невязкам обоих наблюдаемых полей.
